\answer{ex:17:1}{$0.2$~W frente a $81$~\textmu W; el flujo bruto necesitaría como mucho $50$~pJ/bit.}
\answer{ex:17:2}{Con $c=0.1$, $5.6$, $8.7$, $9.2$~bit/s, límite $9.3$~bit/s; con $c=0$ y $N=1000$, $65$~bit/s.}
\answer{ex:17:3}{$B=\log_2N+P\log_2P+(1-P)\log_2\frac{1-P}{N-1}=4.87$, $4.37$, $3.29$~bits por carácter, es decir, $7.3$, $6.6$, $4.9$~bit/s; para $10$~bit/s hacen falta $2.3$~caracteres por segundo.}
\answer{ex:17:4}{Varianza del error $2b/(Nm^2T)+\sigma_\xi^2$ por componente; las contribuciones se igualan con $N\approx90$; límite $\tfrac12\log_2(1+0.25/0.04)\approx1.4$~bits por componente y ventana.}
\answer{ex:17:5}{La pareja converge si $\eta_bd^2+\eta_db^2<2$, es decir, aproximadamente con $\eta<1$: con $0.8$ converge; con $1.1$, oscilaciones sostenidas de periodo $2$; con $1.5$, no periódicas; con $2$, divergencia. Hay que separar las velocidades de los aprendices.}
\answer{ex:17:6}{Umbral $\approx4.4\sigma$; $102$~kbit/s, $0.10$~mW frente a $205$~mW para el flujo bruto, $2000$ veces menos; solo se saturan (más de $3$ espigas) $5.7\cdot10^{-5}$ de las ventanas.}
\answer{ex:17:7}{Ejercicio abierto. Unos $46$~bit/s con $\text{SNR}\approx0.9$ y unos $330$~bit/s con ruido independiente. Si estorba el ruido parecido a la señal, la información extraída se satura al crecer el número de canales; si es el desconocimiento del código, crece con un decodificador mejor (por ejemplo, uno que tenga en cuenta los tiempos de las espigas).}
